\subsection*{一步一步建立直观}

\paragraph{1. 极限关注附近，不强迫取到目标点}

先观察非零步长下的割线。极限问题不是要求在 $h=0$ 计算差商，而是问：当 $h$ 从允许的非零值靠近零时，差商能否被逼近到某个固定数。右图保留两个点，让“接近”与“已经重合”在操作上分开。



\paragraph{2. 误差应当可以被事先控制}

平方函数的差商与 $2a$ 相差恰好 $h$。因此只要要求 $|h|$ 小于给定误差，就能确保读数达到目标精度。这比观察小数位看起来稳定更强：它说明对任意要求的精度，都有一个明确的输入范围保证结果。



\paragraph{3. 检查两侧与定义域}

改变步长的正负，从观察点两侧接近。两侧趋向同一值，才得到通常的双侧导数；若函数只在一侧有定义，应声明研究的是单侧极限。窗口提供的其他函数也有自己的定义域，尤其不能在实对数的非正输入处假装得到函数值。


